% diagrams/firstwords/churros.tex -- ¡Dani se come seis churros con
% azúcar! Uno en la mano, y los otros cinco, en el plato.
\begin{tikzpicture}[dibujofw]
  \newcommand{\churro}[1]{%
    \filldraw[fill=white, rounded corners=1.5mm] (0,-0.22) -- (#1,-0.22) -- (#1,0.22) -- (0,0.22) -- cycle;
    \draw (0.25,0) -- ({#1-0.25},0);}
  \ninoFW{Dani}{Abajo}{Nada}
  \begin{scope}[shift={(0,6.3)}] \bocaOFW \end{scope}
  % el brazo, con el churro a la boca
  \filldraw[fill=white, rounded corners=2mm] (0.75,4.65) -- (1.9,3.7) -- (1.75,5.55) -- (1.45,5.5) -- (1.55,4.25) -- (0.8,4.2) -- cycle;
  \begin{scope}[shift={(0.05,5.75)}] \churro{1.75} \end{scope}
  \filldraw[fill=white] (1.62,5.72) circle (0.27);
  % la mesa y el plato con los otros cinco
  \mesaFW{-3.4}{3.4}{2.2}
  \filldraw[fill=white] (0,2.95) ellipse (2.6 and 0.35);
  \foreach \x/\y/\a in {-1.9/3.1/5, -0.9/3.25/-4, 0.3/3.1/3, -1.4/3.55/-3, 0.2/3.6/6} {\begin{scope}[shift={(\x,\y)}, rotate=\a] \churro{1.6} \end{scope}}
  \draw (-4.2,0) -- (4.2,0);
\end{tikzpicture}
